\documentclass[10pt,a4paper]{amsart}
\usepackage[margin=19mm]{geometry}
\usepackage{amsmath,amssymb,mathtools}
\usepackage[scr=rsfs,cal=euler]{mathalfa}
\usepackage{xfrac}
\usepackage[unicode,hidelinks,bookmarks=false]{hyperref}
\newcommand{\ie}{i.e.\ }
\newcommand{\cf}{cf.\ }
\newcommand{\resp}{respectively\ }
\newcommand{\Subsection}{\subsection}
\providecommand{\nobreakdash}{\nobreak\hbox{-}}
\setlength{\emergencystretch}{2em}
\multlinegap0pt
\usepackage{amssymb}
\newtheorem{theorem}{Theorem}
\newtheorem{lemma}{Lemma}
\newtheorem{claim}{Claim}
\title{Shadowing and Stability of Non-Invertible $p$-adic Dynamics}
\author{D.A. Caprio, F. Lenarduzzi, A. Messaoudi, I. Tsokanos}
\date{Source: arXiv:2408.04779v3 (CC BY 4.0)}
\begin{document}
\maketitle

\noindent\emph{Source and licence.} Shadowing and Stability of Non-Invertible $p$-adic Dynamics. Version arXiv:2408.04779v3 (CC BY 4.0), distributed by arXiv under the Creative Commons Attribution 4.0 International licence, http://creativecommons.org/licenses/by/4.0/, as recorded in the arXiv licence metadata for this exact version and shown on the abstract page https://arxiv.org/abs/2408.04779v3. Full licence notice: \texttt{licenses/arxiv-2408.04779-CC-BY-4.0.txt}.\par\smallskip

\noindent\textit{Editorial scope.} Counted unit: the article's theorem p-adics (source0.tex lines 353--371). The article's complete original proof of it is delivered with it (source0.tex lines 1116--1395, one \texttt{proof} environment of 280 lines, which the article places away from the statement). No mathematics is written, repaired or re-derived: the statement and the proof are the article's own lines, in the article's own notation, and no retained line is substituted or re-flowed.  Retained uncounted as the source-local context the proof needs: lines 969--974; lines 1038--1042; lines 1072--1079; lines 1075--1077.  Nothing was omitted from any retained span, so the package carries no omission record.  No retained line contains a citation, so the wrapper carries no bibliography and no bibliography entry.  No retained line contains a figure environment, an image-inclusion command or drawing code, so no image is delivered and the package declares no asset dependency.  Editorial formatting only, in addition: an AMS \texttt{amsart} preamble that restates the article's own declarations and macros used by the retained lines, namely the environments \texttt{theorem}, \texttt{lemma}, \texttt{claim} on separate counters, together with the standard packages those lines need.  The retained environments print in this document as Theorem 1, Lemma 1, Claim 1 in order of appearance, whereas the article prints the same items with the numbering of its own full document.  The complete native source of the article as distributed in the arXiv e-print for this exact version is bundled unchanged as \texttt{sources/163549/source0.tex}.
	\begin{lemma}\label{Lem2_Open bi-Lipschitz Contractions}
		Let $R: \mathbb{X}_{p} \to \mathbb{X}_{p}$, with $\mathbb{X}_{p} \in \{\mathbb{Z}_{p}, \mathbb{Q}_{p}\}$, be a contraction satisfying the bi-Lipschitz inequality \eqref{Eq2_Bi-Lipschitz Contractions}. 
		If the image $R(\mathbb{X}_{p})$ is open, then $R$ is an open mapping. 
		In particular, if $R$ is scaling, then its image $R(\mathbb{X}_{p})$ is $\rho$-open for some $\rho > 0$.
		
	\end{lemma}

	\begin{lemma}\label{Lem2_Open bi-Lipschitz Contraction plus Lipschitz}
		Let $R: \mathbb{X}_{p} \to \mathbb{X}_{p}$, with $\mathbb{X}_{p} \in  \left\{  \mathbb{Z}_{p}, \mathbb{Q}_{p}  \right\} $, be a contraction satisfying the bi-Lipschitz inequality \eqref{Eq2_Bi-Lipschitz Contractions} with constants $0  < c_{1} \le c_{2}  < 1$. 
		Then, for every $0  < \delta \le  c_{1}/2  $ and every $\phi \in Lip_{\delta} \left(  \mathbb{X}_{p}  \right) $, the map $T = R + \phi$ is a bi-Lipschitz contraction satisfying the bi-Lipschitz inequality \eqref{Eq2_Bi-Lipschitz Contractions} with the same constants $c_{1}$ and $c_{2}$. 
		Moreover, if the image $R \left(  \mathbb{X}_{p}  \right) $ is a $\rho$-open set for some $\rho  > 0$ and $0  < \delta \le \min \left\{  c_{1}/2,   \rho/2  \right\} $, then $ T \left(  \mathbb{X}_{p}  \right)  = R \left(  \mathbb{X}_{p}  \right) $.
	\end{lemma}

	\begin{lemma}\label{Lem2_Scaling Contraction plus Lipschitz}
		Let $R: \mathbb{X}_{p} \to \mathbb{X}_{p}$, with $\mathbb{X}_{p} \in  \left\{  \mathbb{Z}_{p}, \mathbb{Q}_{p}  \right\} $, be a contraction satisfying the left-hand side of the bi-Lipschitz inequality \eqref{Eq2_Bi-Lipschitz Contractions} with constant $0  < c_{1}   < 1$. 
		If $R$ is a scaling contraction, then for every $0  < \delta \le  c_{1}/2  $ and for every $\phi \in Lip_{\delta} \left(  \mathbb{X}_{p}  \right) $, the map $T = R + \phi$ is scaling and satisfies 
		\begin{equation}\label{EqLem_Scaling Maps Induction}
			\left| \left|  T^{n}(x) - T^{n}(y)  \right|  \right|_{p}   =    \left| \left|  R^{n}(x) - R^{n}(y)  \right|  \right|_{p}, \quad \text{for every } x,y \in \mathbb{X}_{p} \text{ and } n \in \mathbb{N} . 
		\end{equation} 
		
	\end{lemma}

		\begin{equation}
			\left| \left|  T^{n}(x) - T^{n}(y)  \right|  \right|_{p}   =    \left| \left|  R^{n}(x) - R^{n}(y)  \right|  \right|_{p}, \quad \text{for every } x,y \in \mathbb{X}_{p} \text{ and } n \in \mathbb{N} . 
		\end{equation} 

	\begin{theorem}\label{Theor3_Contractions in p-adics}  
		Let $R: \mathbb{X}_{p} \to \mathbb{X}_{p}$ be a contraction whose image $R \left(  \mathbb{X}_{p}  \right) $ is open and which satisfies the bi-Lipschitz inequality 
		\begin{equation}\label{Eq2_Bi-Lipschitz Contractions}
			c_{1}\cdot  \left| \left|  x -y  \right|  \right|_{p}   \le       \left| \left|  R(x) - R(y)  \right|  \right|_{p}   \le   c_{2} \cdot  \left| \left|  x -y  \right|  \right|_{p}, \quad \text{for every }  x,y \in \mathbb{X}_{p},  
		\end{equation}
		where $ 0  < c_{1} \le c_{2}  < 1 $ are positive constants.  Then, the following statements hold:
		\begin{enumerate}
			
			
			\item Assume that $\mathbb{X}_{p} = \mathbb{Z}_{p}$. Then $R$ is strongly Lipschitz structurally stable. 
			
			
			\item Assume that $\mathbb{X}_{p} = \mathbb{Q}_{p}$. If $R$ is scaling, then $R$ is strongly Lipschitz structurally stable.
			
			
			
		\end{enumerate}
		
	\end{theorem}

	\begin{proof}[Proof of Theorem \ref{Theor3_Contractions in p-adics}] 
		
		Fix $\mathbb{X}_{p} \in  \left\{  \mathbb{Z}_{p}, \mathbb{Q}_{p}  \right\} $ and let $R: \mathbb{X}_{p} \to \mathbb{X}_{p}$ be a bi-Lipschitz contraction on $\mathbb{X}_{p}$ that satisfies the bi-Lipschitz inequality \eqref{Eq2_Bi-Lipschitz Contractions} with positive constants $0  < c_{1} \le c_{2}  < 1$. Let $x_{R} \in \mathbb{X}_{p}$ denote the unique fixed point of $R$. 
		
		
		
		\paragraph{}
		
		Assume, without loss of generality, that $R \left(  \mathbb{X}_{p}  \right) $ is $\rho$-open for some $\rho > 0$. 
		This assumption is satisfied in the present setting. 
		In the case of Point~(1) of Theorem~\ref{Theor3_Contractions in p-adics}, the condition holds because $R \left(  \mathbb{Z}_{p}  \right) $ is open by hypothesis and compact by continuity. 
		For Point~(2), where $R$ is assumed to be scaling, Lemma~\ref{Lem2_Open bi-Lipschitz Contractions} ensures that the image $R \left(  \mathbb{X}_{p}  \right) $ is $\rho$-open.
		
		
		\paragraph{}
		
		We now present a unified argument covering both Point~(1) and Point~(2). The specific assumptions required to complete the proof in each case are invoked only at the final stage, in the respective concluding paragraphs. 
		
		
		
		\paragraph{}
		
		
		
		Fix $0  < \delta  <  \min  \left\{  \rho/2, c_{1}/2  \right\} $, $\phi \in Lip_{\delta} \left(  \mathbb{X}_{p}  \right) $, and set $T = R + \phi$. By Lemma \ref{Lem2_Open bi-Lipschitz Contraction plus Lipschitz}, the map $T$ is also a bi-Lipschitz contraction satisfying   
		\begin{equation}\label{Eq3_Contraction T Lower Bound}
			c_{1} \cdot  \left| \left|  x - y  \right|  \right|_{p}   \le    \left| \left|  T(x) - T(y)  \right|  \right|_{p}    \le   c_{2} \cdot  \left| \left|  x - y  \right|  \right|_{p}, \quad \text{for all } x,y \in \mathbb{X}_{p}, 
		\end{equation}
		where the constants $c_{1}$ and $ c_{2} $ are the same as in the inequality \eqref{Eq2_Bi-Lipschitz Contractions}. Denote by $x_{T} \in \mathbb{X}_{p}$ the unique fixed point of $T$. 
		
		Lemmas \ref{Lem2_Open bi-Lipschitz Contractions} and \ref{Lem2_Open bi-Lipschitz Contraction plus Lipschitz} together imply that both $R$ and $T$ are open maps with  $ T \left(  \mathbb{X}_{p}  \right)  = R \left(  \mathbb{X}_{p}  \right)  $. Thus, one can define the set  
		\begin{equation}\label{Eq2_Domain of the Contraction}  
			\mathcal{V} :=    \mathbb{X}_{p} \backslash R \left(  \mathbb{X}_{p}  \right)    =   \mathbb{X}_{p} \backslash T \left(  \mathbb{X}_{p}  \right)  .
		\end{equation}
		Note that $\mathcal{V}$ is a $\rho$-open set, as it is the complement of the $\rho$-open set $R \left(  \mathbb{X}_{p}  \right) $ in $\mathbb{X}_{p}$. 
		We also introduce the sets 
		\begin{equation}\label{Eq2_Contraction Domain B}
			A_{R} :=   \bigcup_{n=0}^{+ \infty } R^{n} \left(  \mathcal{V}  \right) , \quad B_{R} :=   \mathbb{X}_{p} \backslash A_{R}, \quad A_{T} :=   \bigcup_{n=0}^{+ \infty } T^{n} \left(  \mathcal{V}  \right) ,   \text{and }   B_{T} :=   \mathbb{X}_{p} \backslash A_{T}.
		\end{equation}
		
		
		\paragraph{}
		
		The proof will now proceed by constructing a homeomorphism $h: \mathbb{X}_{p} \to \mathbb{X}_{p}$ such that $R \circ h = h \circ T$.  Specifically, $h$ is given by the formula 
		\begin{equation}\label{Eq3_Conjugation Fundamental Domains A and B}
			h \left(  x  \right)  :=   \begin{cases}
				h_{1} \left(  x  \right) ,  & \text{if } x \in A_{T},  \\ 
				h_{2} \left(  x  \right) ,   & \text{if } x \in B_{T},  \\  
			\end{cases}
		\end{equation}
		where $h_{1}: A_{T} \to A_{R}$ and $h_{2}: B_{T} \to B_{R}$ are bijections, which will be specified below. Along with the definitions of $h_{1}$ and $h_{2}$, we will show that 
		\begin{equation}\label{Eq3_Conjugation h delta close to Identity}
			\left| \left|  h_{1}(x) - x  \right|  \right|_{p}   \le   \delta,   \text{ for all } x \in A_{T}, \text{ and }    \left| \left|  h_{2}(x) - x  \right|  \right|_{p}   \le   \delta,   \text{ for all } x \in B_{T}. 
		\end{equation}
		Finally, we will prove the continuity of $h$ and $h^{-1}$. 
		
		
		To define the maps $h_{1}$ and $h_{2}$, we use the following lemma, which can be readily proved.   
		
		\begin{lemma}\label{Lem2_Domain of the Contraction}
			Let $R: \mathbb{X}_{p} \to \mathbb{X}_{p}$ be a left invertible contraction. Then 
			\begin{equation}\label{Eq2_Partition Domain}
				\mathbb{X}_{p}   =   B_{R} \cup \bigcup_{n=0}^{+ \infty } R^{n} \left(  \mathcal{V}  \right)  , 
			\end{equation}
			where $ \mathcal{V} $ and $B_{R}$ are as defined in relations \eqref{Eq2_Domain of the Contraction} and \eqref{Eq2_Contraction Domain B}, respectively.
			In particular, the sets $B_{R}$ and $R^{n} \left(  \mathcal{V}  \right) $, for $n \in \mathbb{N} $,  partition $\mathbb{X}_{p}$ and satisfy 
			\begin{equation}\label{Eq2_Partition Domain Explicit}
				R^{n} \left(  \mathcal{V}  \right)    =   R^{n} \left(  \mathbb{X}_{p}  \right)  \backslash R^{n+1} \left(  \mathbb{X}_{p}  \right) , \quad \text{for every }  n \in  \mathbb{N} .  
			\end{equation}
			
			Moreover, the restriction of $R$ to $B_{R}$, i.e., $R: B_{R} \to B_{R}$, is a homeomorphism, and $x_{R} \in B_{R}$. 
			
		\end{lemma}
		
		\paragraph{}
		
		To define the bijection $h_{1}$, observe that inequalities \eqref{Eq2_Bi-Lipschitz Contractions} and \eqref{Eq3_Contraction T Lower Bound} imply the existence of two continuous bijections $L: R \left(  \mathbb{X}_{p}  \right)  \to \mathbb{X}_{p}$ and $\Tilde{L}: T \left(  \mathbb{X}_{p}  \right)  \to \mathbb{X}_{p}$ that satisfy the relations $L \circ R = id$ and $ \Tilde{L} \circ T = id$, respectively. Therefore, Lemma \ref{Lem2_Domain of the Contraction} guarantees that the map $h_{1}: A_{T} \to A_{R}$, defined by the formula  
		\begin{equation}\label{Eq3_Conjugation Fundamental Domain A}
			h_{1} \left(  x  \right)   :=   R^{n} \left(  v  \right) , \quad \text{where } x = T^{n} \left(  v  \right)  \in A_{T}   \text{ for some } v \in \mathcal{V} \text{ and } n \in  \mathbb{N} ,   
		\end{equation} 
		is a well-defined bijection. 
		Moreover, $h_{1}$  satisfies the left-hand side of \eqref{Eq3_Conjugation h delta close to Identity}. Indeed, fix $ x \in A_{T}$, where $x = T^{n} \left(  v  \right) $ for some $v \in \mathcal{V}$ and $n \in \mathbb{N} $. Then, one has that $h_{1}(x) - x = R^{n}(v) - T^{n}(v)$. 
		If $n=0$, it is clear that $ \left| \left|  R^{0}(v) - T^{0}(v)  \right|  \right|_{p} = 0$. If $n \ge 1$, then the claim follows inductively from the inequality 
		\begin{align*}
			\left| \left|  R^{n}(v) - T^{n}(v)  \right|  \right|_{p}   & =    \left| \left|  R\circ R^{n-1} \left(  v  \right)  - R\circ T^{n-1} \left(  v  \right)  - \phi\circ T^{n-1} \left(  v  \right)   \right|  \right|_{p} \\ 
			& \underset{\eqref{Eq2_Bi-Lipschitz Contractions}}{\le}   \max \left\{  \delta, c_{2}\cdot  \left| \left|  R^{n-1}(v) - T^{n-1}(v)  \right|  \right|_{p}  \right\}  .
		\end{align*}
		Finally, observe that the map $h_{1}$ conjugates $R$ and $T$. Indeed, one has that  
		\begin{equation}\label{Eq3_Homeomorphism 1 Conjugation}
			R\circ h_{1}(x)    \underset{\eqref{Eq3_Conjugation Fundamental Domain A}}{=}    R^{n+1}(v)    \underset{\eqref{Eq3_Conjugation Fundamental Domain A}}{=}    h_{1}\circ T^{n+1}(v)    \underset{\eqref{Eq3_Conjugation Fundamental Domain A}}{=}    h_{1}\circ T(x) .
		\end{equation}
		
		
		
		\paragraph{} 
		
		To define the bijection $h_{2} : B_{T} \to B_{R} $, let 
		$$  \ell^{ \infty } \left(  \mathbb{Z} , \mathbb{X}_{p}  \right)  : =    \left\{   \left(  x_{n}   \right)_{n \in \mathbb{Z} } \in \mathbb{X}_{p}^{\mathbb{Z} } :   \sup_{n \in \mathbb{Z} }  \left| \left|  x_{n}  \right|  \right|_{p}   <   + \infty   \right\}  $$
		be the Banach space of bi-infinite bounded sequences in $\mathbb{X}_{p}$. 
		Endowed with the supremum norm $ \left| \left|  \cdot  \right|  \right|_{\ell^{ \infty } \left(  \mathbb{Z} , \mathbb{X}_{p}  \right) }$, the space $ \left(  \ell^{ \infty } \left(  \mathbb{Z} , \mathbb{X}_{p}  \right) ,  \left| \left|  \cdot  \right|  \right|_{\ell^{ \infty } \left(  \mathbb{Z} , \mathbb{X}_{p}  \right) }  \right) $ is complete. Notice that $\ell^{ \infty } \left(  \mathbb{Z} , \mathbb{Z}_{p}  \right) $ is the closed unit ball of the $\mathbb{Q}_{p}$-Banach space $\ell^{ \infty } \left(  \mathbb{Z} , \mathbb{Q}_{p}  \right) $ and, in particular, $\ell^{ \infty } \left(  \mathbb{Z} , \mathbb{Z}_{p}  \right)  = \mathbb{Z}_{p}^{\mathbb{Z} }$. 
		
		Fix $x \in B_{T}$. By Lemma \ref{Lem2_Domain of the Contraction}, the restriction $T: B_{T} \to B_{T}$ is a homeomorphism. Consequently, for all $x \in B_{T}$, the iterates $T^{n} \left(  x  \right)  \in B_{T}$ are well defined for all $n \in \mathbb{Z} $. This allows the definition of the map $W_{x}: \ell^{ \infty } \left(  \mathbb{Z} , \mathbb{X}_{p}  \right)  \to \ell^{ \infty } \left(  \mathbb{Z} , \mathbb{X}_{p}  \right) $ by  
		\begin{equation}\label{Eq3_Fixed Point Operator}
			W_{x} \left(   \left(  u_{n}   \right)_{n \in \mathbb{Z} }  \right)  :=    \left(  R \left(  T^{n-1}(x) + u_{n-1}   \right)  -  T^{n}(x)      \right)_{n \in \mathbb{Z} } .
		\end{equation} 
		This map is well-defined and is a contraction on $\ell^{ \infty } \left(  \mathbb{Z} , \mathbb{X}_{p}  \right) $, and therefore admits a unique fixed point $ \left(  z_{n} \left(  x  \right)    \right)_{n \in \mathbb{Z} }$, which in particular satisfies 
		\begin{equation}\label{Eq3_Fixed Point Recursive Relation} 
			R \left(  T^{n-1} \left(  x  \right)  + z_{n-1} \left(  x  \right)   \right)    =   T^{n} \left(  x  \right)  + z_{n} \left(  x  \right)  , \quad \text{for all } n \in \mathbb{Z} .
		\end{equation}
		
		
		
		Define $h_{2}: B_{T} \to B_{R}$ by  
		\begin{equation}\label{Eq2_Conjugation Fundamental Domain B}
			h_{2} \left(  x  \right)  :=   x + z_{0} \left(  x  \right)  
		\end{equation}
		To verify that $h_{2}$ is well defined, it remains to show that $h_{2} \left(  x  \right)  \in B_{R}$. To this end, relation \eqref{Eq3_Fixed Point Recursive Relation} implies that 
		\begin{equation}\label{Eq3_Inverse Images of B_R Points}
			T^{-n} \left(  x  \right)  + z_{-n} \left(  x  \right)    \in   R^{-n} \left\{  h_{2}  \left(  x  \right)   \right\}  \quad \text{for every } n \ge 0 .
		\end{equation} 
		This yields the claim immediately. Indeed, if $ h_{2}(x) \in A_{R}$, then there exist $v \in \mathcal{V}$ and $n_{0} \in \mathbb{N} $ such that $ h_{2}(x) = R^{n_{0}}(v)$. 
		In that case, $R^{-n_{0} -1 } \left\{   h_{2}(x)  \right\}  = \emptyset$, which contradicts~\eqref{Eq3_Inverse Images of B_R Points}. 
		
		
		
		
		The map $h_{2}: B_{T} \to B_{R}$ is a bijection. To prove this, we define the map $\tilde{h}_{2}: B_{R} \to B_{T}$ in the same way as $h_{2}$, but with the roles of $R$ and $T$ reversed. 
		For every $x \in B_{R}$ and $y \in B_{T}$, the uniqueness of the points $h_{2} \left(  x  \right) $ and $\tilde{h}_{2} \left(  y  \right) $ implies that  $h_{2} \circ \tilde{h}_{2} \left(  y  \right)  = y$ and $\tilde{h}_{2} \circ h_{2} \left(  x  \right)   = x$. Thus, we conclude that $h_{2}$ is a bijection. 
		
		
		To demonstrate that $h_{2}$ satisfies the right-hand side inequality of \eqref{Eq3_Conjugation h delta close to Identity}, we begin by observing the following:   
		\begin{align*}  \left| \left|  z_{n} \left(  x  \right)   \right|  \right|_{p}   &
			\underset{\eqref{Eq3_Fixed Point Operator}}{=}    \left| \left|  R \left(  T^{n-1} \left(  x  \right)  + z_{n-1} \left(  x  \right)   \right)  - T^{n} \left(  x  \right)    \right|  \right|_{p}  \\  
			& =    \left| \left|  R \left(  T^{n-1} \left(  x  \right)  + z_{n-1} \left(  x  \right)   \right)  - R \left(  T^{n-1} \left(  x  \right)   \right)  - \phi \left(  T^{n-1} \left(  x  \right)   \right)    \right|  \right|_{p}   \\  
			& \le    \max \left\{  c_{2}  \left| \left|  z_{n-1} \left(  x  \right)   \right|  \right|_{p},   \delta  \right\} ,
		\end{align*}
		for all $n \in \mathbb{Z} $. Since $ \left(  z_{n}   \right)_{n \in \mathbb{Z} } \in \ell^{ \infty } \left(  \mathbb{Z} , \mathbb{X}_{p}  \right) $, this recurrence relation implies that:  
		\begin{equation}\label{Eq3_Fixed Point delta small}
			\left| \left|  z_{n} \left(  x  \right)   \right|  \right|_{p}   \le   \delta, \quad \text{for every } x \in B_{T} \text{ and } n \in \mathbb{Z} . 
		\end{equation}
		This inequality ensures that $h_{2}$ satisfies the desired bound in \eqref{Eq3_Conjugation h delta close to Identity}.
		
		
		Finally, observe that the map $h_{2}$ conjugates $R$ and $T$. 
		Indeed, since the maps $W_{x}, W_{T \left(  x  \right) } : \ell^{ \infty }_{\mathbb{Z} } \left(  \mathbb{X}_{p}  \right)  \to \ell^{ \infty }_{\mathbb{Z} } \left(  \mathbb{X}_{p}  \right) $ are contractions, each admits a unique fixed point. Applying relation~\eqref{Eq3_Fixed Point Recursive Relation} to $x$ and $T \left(  x  \right) $, respectively, yields $ z_{n} \left(  x  \right)  = z_{n - 1} \left(  T \left(  x  \right)   \right) $, for all $n \in \mathbb{Z} $. Consequently,  
		\begin{align*} 
			R \circ h_{2} \left(  x  \right)    &\underset{\eqref{Eq2_Conjugation Fundamental Domain B}}{=}   R \left(  x + z_{0} \left(  x  \right)   \right)    \underset{\eqref{Eq3_Fixed Point Recursive Relation}}{=}   T \left(  x  \right)  + z_{1} \left(  x  \right)   \\  
			&=   T \left(  x  \right)  + z_{0} \left(  T \left(  x  \right)   \right)    =   h_{2} \circ T \left(  x  \right)  .
		\end{align*}
		Therefore, for every $x \in B_{T}$,  
		\begin{equation}\label{Eq3_Homeorphism 2 Conjugation}
			R\circ h_{2} \left(  x  \right)    =   h_{2} \circ T(x) .
		\end{equation}
		
		
		
		\paragraph{} 
		
		To prove that $R$ is strongly Lipschitz structurally stable, it remains to show that the conjugation map $h$ and its inverse $h^{-1}$ are continuous. 
		We prove the result for $h$ only, as the proof for $h^{-1}$ follows similarly. 
		
		\medskip
		
		First, it is shown that $h$ is continuous at every point $x \in A_{T}$. The set $A_{T}$ is open, being the union of the sets $T^{n} \left(  \mathcal{V}  \right) $ for $n \in \mathbb{N} $, where each $T^{n} \left(  \mathcal{V}  \right)  $ is open because $T$ is an open map and $\mathcal{V}$ is $\rho$-open. 
		Fix $x \in A_{T} $, and let $m \in \mathbb{N} $ and $v \in \mathcal{V}$ be such that $x = T^{m}(v)$. Choose $  \epsilon > 0$ sufficiently small so that the radius 
		$$  \eta :=   {c_{1}^{m} \over c_{2}^{m}} \cdot  \epsilon $$
		satisfies $B \left(  x,  \eta  \right)  \subseteq T^{m} \left(  \mathcal{V}  \right)  $. 
		For every $x' \in B \left(  x ,  \eta  \right) $, there exists $v' \in \mathcal{V}$ such that $x' = T^{m} \left(  v'  \right) $. Inequality~\eqref{Eq3_Contraction T Lower Bound} then yields  $ \left| \left|  v - v'  \right|  \right|_{p} \le \epsilon / c_{2}^{m}$, while inequality~\eqref{Eq2_Bi-Lipschitz Contractions} implies  
		$$  \left| \left|  h \left(  x  \right)  - h \left(  x'  \right)   \right|  \right|_{p}   =    \left| \left|  h_{1}(x) - h_{1} \left(  x' \right)   \right|  \right|_{p}   =    \left| \left|  R^{m} \left(  v  \right)  - R^{m} \left(  v'  \right)   \right|  \right|_{p}   \le    \epsilon . $$ 	
		Since $x \in A_{T}$ and $\epsilon  > 0$ were arbitrary, this proves that $h$ is continuous on $A_{T}$. 
		
		
		\medskip
		
		The continuity of $h$ at points $x \in B_{T}$ is established separately in each of the cases considered in the theorem. More precisely, in the setting of Point~(1) the argument relies on the compactness of the space $\mathbb{Z}_{p}$, whereas in the setting of Point~(2) it uses the assumption that $R$ is scaling.
		
		
		\paragraph{Proof in the case of Point (1).} 
		
		Assume that $\mathbb{X}_{p} = \mathbb{Z}_{p}$. 
		Note that $B_{T} =  \left\{  x_{T}  \right\} $ and $B_{R} =  \left\{  x_{R}  \right\} $. To demonstrate this, assume there exists $x \in B_{T}$ such that $x \neq x_{T}$. By the construction of the set $B_{T}$, one has that $T \left(  B_{T}  \right)  = B_{T}$. Let $\textrm{diam} \left(  B_{T}  \right) $ be the diameter of $ B_{T}$, defined as   
		$$ \textrm{diam} \left(  B_{T}  \right)  :=   \sup \left\{   \left| \left|  x_{1} - x_{2}  \right|  \right|_{p} :   x_{1}, x_{2} \in B_{T}  \right\}    <   + \infty  .$$
		Assume $\textrm{diam} \left(  B_{T}  \right)   > 0$. Since $R$ is a contraction, one has that 
		$$ \textrm{diam} \left(  B_{T}  \right)    =   \textrm{diam} \left(  T \left(  B_{T}  \right)   \right)     \le   c_{2} \cdot \textrm{diam} \left(  B_{T}  \right)  ,$$
		which contradicts the assumption that $\textrm{diam} \left(  B_{T}  \right)   > 0$. Therefore, $\textrm{diam} \left(  B_{T}  \right)  =0$, and thus $ B_{T}  =  \left\{  x_{T}  \right\} $. 
		A similar argument applies for $B_{R}$, yielding $B_{R} =  \left\{  x_{R}  \right\} $. 
		
		
		To conclude the proof, it remains to show that $h$ is continuous at $x_{T}$. 
		To this end, fix $ \eta  > 0$ and assume that $x \in \mathbb{Z}_{p}$ satisfies $ \left| \left|  x_{T} - x  \right|  \right|_{p} \le  \eta $ and $x \neq x_{T}$. Let $v = v \left(  x  \right)  \in \mathcal{V} $ and $n = n \left(  x  \right)  \in \mathbb{N} $ be such that $x = T^{n}(v)$. 
		It suffices to prove that $n \left(  x  \right)  \to + \infty $ as $x \to x_{T}$, since
		\begin{equation}\label{Eq2_Fixed Point Continuity}
			\left| \left|  h \left(  x_{T}  \right)  - h \left(  x  \right)   \right|  \right|_{p}   =    \left| \left|  x_{R} - R^{n}(v)  \right|  \right|_{p}    =    \left| \left|  R^{n} \left(  x_{R}  \right)  - R^{n}(v)  \right|  \right|_{p}   \underset{\eqref{Eq2_Bi-Lipschitz Contractions}}{\le}   c_{2}^{n} . 
		\end{equation}
		To this end, for every $k \in \mathbb{N} $, define $\eta_{k} := \inf  \left\{   \left| \left|  x_{T} - x'  \right|  \right|_{p} :   x' \in T^{k} \left(  \mathcal{V}  \right)   \right\} $.
		Due to the compactness of $\mathcal{V}$, one infers that $\eta_{k}  > 0$ for all $k \in \mathbb{N} $.
		Moreover, since $T$ is a contraction, the sequence $ \left\{  \eta_{k}  \right\}_{k \in  \mathbb{N} }$ is strictly decreasing, and thus $\eta_{k} {\longrightarrow} 0$ as $k \to + \infty $. 
		Let $k_{ \eta} \in \mathbb{N} $ be such that $\eta_{k_{ \eta} + 1}  <  \eta \le \eta_{k_{ \eta}} $. From the definition of the sequence $ \left(  \eta_{k}   \right)_{k \in  \mathbb{N} }$, one has that $n \ge k \left(   \eta  \right)  $. 
		Therefore, as $ \eta \to 0$, one obtains $\underset{ \eta \to 0}{\lim} k_{ \eta} = + \infty $, which implies that $\underset{ x \to x_{T} }{\lim} n(x) = + \infty $. 
		This shows that $h$ is continuous at $x_{T}$. 
		
		
		
		
		\paragraph{Proof in the case of Point (2).} Assume that $\mathbb{X}_{p} = \mathbb{Q}_{p}$ and that $R$ is scaling. Recall that, by Lemma \ref{Lem2_Open bi-Lipschitz Contractions}, the map $R$ is open and its image $R \left(  \mathbb{X}_{p}  \right) $ is $\rho$-open for some $\rho  > 0$. 
		
		
		
		
		
		\paragraph{} 	
		
		To prove that $h$ is continuous on $B_{T}$, we will use two claims. In the following, let $x \in B_{T}$ be fixed, and we will establish the continuity of $h$ at this point.
		
		\begin{claim}\label{Claim3 Homeomorphism 2 on B_T}
			For every $ \epsilon  > 0$, there exists $ 0 <  \eta <  \epsilon $ such that for all $x' \in B_{T}$ with $ \left| \left|  x - x'  \right|  \right|_{p} \le  \eta $, we have $  \left| \left|  h_{2} \left(  x  \right)   - h_{2} \left(  x'  \right)   \right|  \right|_{p}  \le  \epsilon $. 
			
		\end{claim}
		
		\begin{proof}[Proof of Claim \ref{Claim3 Homeomorphism 2 on B_T}]
			
			Fix $  \epsilon  > 0$ and $x' \in B_{T}$ such that $ \left| \left|  x - x'  \right|  \right|_{p} \le  \eta $, where $ 0 <  \eta <  \epsilon $ is a positive number that will be determined later. 
			To prove that $ \left| \left|  h_{2}(x) - h_{2} \left(  x'  \right)   \right|  \right|_{p} \le  \epsilon $, it suffices to show that $ \left| \left|  z_{0}(x) - z_{0} \left(  x'  \right)   \right|  \right|_{p} \le  \epsilon $, based on relation \eqref{Eq2_Conjugation Fundamental Domain B}.    
			Since $ \left(  z_{n} \left(  x  \right)    \right)_{n \in \mathbb{Z} } \in \ell^{ \infty } \left(  \mathbb{Z} , \mathbb{X}_{p}  \right) $ and $ \left(  z_{n} \left(  x'  \right)    \right)_{n \in \mathbb{Z} } \in \ell^{ \infty } \left(  \mathbb{Z} , \mathbb{X}_{p}  \right) $ are the fixed point of the maps $W_{x}$ and $W_{x'}$, respectively, from relation \eqref{Eq3_Fixed Point Operator}, one infers that for every $m \in \mathbb{Z} $ it holds that
			
			\begin{align*} 
				\left| \left|  z_{m} \left(  x  \right)  - z_{m} \left(  x'  \right)   \right|  \right|_{p}   \le    \max \left\{  \begin{matrix}  \left| \left|  T^{m} \left(  x  \right)  - T^{m} \left(  x'  \right)   \right|  \right|_{p},  \\    c_{2}   \left| \left|  T^{m-1} \left(  x  \right)  - T^{m-1} \left(  x'  \right)   \right|  \right|_{p},  \\ 
					c_{2}  \left| \left|  z_{m-1} \left(  x  \right)  - z_{m-1} \left(  x'  \right)   \right|  \right|_{p} \end{matrix}  \right\}  .
			\end{align*}
			
			By applying this inequality recursively for $m$ ranging from $0$ to $-n$, where $n \ge 0$ is a nonnegative integer, we obtain the following estimate: 
			\begin{align*}
				\left| \left|  z_{0} \left(  x  \right)  - z_{0} \left(  x'  \right)   \right|  \right|_{p}   & \le    \max \left\{  \begin{matrix}   \left| \left|  x - x'  \right|  \right|_{p} ,  \\  
					\max \left\{   c_{2}^{m}  \left| \left|  T^{-m} \left(  x  \right)  - T^{-m} \left(  x'  \right)   \right|  \right|_{p} :   1 \le m \le n  \right\}  ,   \\  
					c_{2}^{n}  \left| \left|  z_{-n} \left(  x  \right)  - z_{-n} \left(  x'  \right)   \right|  \right|_{p} \end{matrix}  \right\}    \\ 
			\end{align*}
			Upon taking into account inequality \eqref{Eq3_Fixed Point delta small}, the right-hand side of this inequality can be made smaller than $ \epsilon  > 0$. This can be done by first selecting $n \in \mathbb{N} $ sufficiently large and then choosing $ \eta  > 0$ small enough. 
			
			The proof of the claim is complete. 
			
			
		\end{proof}
		
		
		\begin{claim}\label{Claim3 Homeomorphism 2 on A_T}
			Let $ \eta  > 0$ be a positive number, and  $T^{n} \left(  v  \right)  \in A_{T}$ for some $v \in \mathcal{V}$  be such that $ \left| \left|  x - T^{n} \left(  v  \right)   \right|  \right|_{p} \le  \eta $. Then, the following inequality holds:  
			$$   \left| \left|  h \left(  x  \right)   - h \left(  T^{n} \left(  v  \right)   \right)   \right|  \right|_{p}    \le     \max \left\{  \delta  c_{2}^{n_{ \eta}} ,    \eta  \right\}  ,$$
			where $n_{ \eta} \in \mathbb{N} $ is the largest non-negative integer such that $   \eta \cdot  \left(  1/ c_{1}  \right) ^{n_{ \eta}} \le \rho $. 
			
		\end{claim}
		
		\begin{proof}[Proof of Claim \ref{Claim3 Homeomorphism 2 on A_T}]
			Fix $  \eta  > 0$ and a point $T^{n} \left(  v  \right)  \in A_{T}$, where $v \in \mathcal{V}$ and $ \left| \left|  x - T^{n}(v)  \right|  \right|_{p} \le  \eta $. The choice of $n_{ \eta} $ guarantees that $n \ge n_{ \eta}$. Indeed, it ensures that  
			$$ \left| \left|  T^{-n_{ \eta}}(x) -  T^{- n_{ \eta}} \left(  T^{n}(v)  \right)   \right|  \right|_{p} \le \rho. $$
			Therefore since $T \left(  \mathbb{Q}_{p}  \right) $ is $\rho$-open, it follows that  $T^{n - n_{ \eta}}(v) \in T \left(  \mathbb{Q}_{p}  \right) $. In other words, the point $T^{n} \left(  v  \right) $ is at least $n_{ \eta}$-times invertible by the map $T$. 
			
			
			By Lemma \ref{Lem2_Scaling Contraction plus Lipschitz}, one has that
			$$  \left| \left|  R^{n} \left(  T^{-n}(x)  \right)  - R^{n}(v)  \right|  \right|_{p}   \underset{\eqref{EqLem_Scaling Maps Induction}}{=}    \left| \left|  x - T^{n}(v)  \right|  \right|_{p} . $$ 
			Furthermore, it holds that: $  \left| \left|  h\circ T^{-n}(x)  - T^{-n}(x)  \right|  \right|_{p} \le \delta$ and $ R^{-n}\circ h(x)  = h\circ T^{-n} (x) $.
			Therefore,  one has that
			\begin{align*}
				\left| \left|  h(x) - h \left(  T^{n}(v)  \right)   \right|  \right|_{p}  & =    \left| \left|  R^{n}\circ R^{-n} \circ h(x) - R^{n}(v)  \right|  \right|_{p}  \\ 
				&=    \left| \left|  R^{n}\circ h \circ T^{-n}(x) - R^{n}\circ T^{-n} (x) + R^{n} \circ T^{-n}(x)   - R^{n}(v)  \right|  \right|_{p}  \\ 
				& \underset{\eqref{EqLem_Scaling Maps Induction}}{\le}   \max  \left\{  \delta  c_{2}^{n},    \left| \left|  x - T^{n}(v)  \right|  \right|_{p}  \right\}     \le   \max \left\{  \delta  c_{2}^{n_{ \eta}} ,    \eta  \right\} , 
			\end{align*}
			where the last inequality follows from the assumption. 
			This completes the proof of the claim. 
			
		\end{proof}
		
		
		From the Claims \ref{Claim3 Homeomorphism 2 on B_T} and \ref{Claim3 Homeomorphism 2 on A_T}, we have shown that $h$ is continuous at every point $x \in B_{T}$. Specifically, the first claim establishes a continuity condition on $h$ for small perturbations within $B_{T}$, while the second claim provides a bound on the continuity of $h$ when $x$ is close to a point in $A_{T}$. Therefore, $h$ is continuous on $B_{T}$. 
		
		
		
		The proof of the theorem is complete. 
		
		
	\end{proof}

\end{document}
